% ============================================================
% setup/fonts.tex —— 字体设置
%
% 要求：英文字体 = Times New Roman（新罗马）
% 中文：沿用 ctex 根据操作系统自动选择的字体
%       （Windows: 宋体/黑体/楷体/仿宋；macOS: 华文系列；Linux: Fandol）
%
% 说明：\setmainfont 只改变"西文（英文）衬线字体"，
%       中文字体仍由 ctex 的 \setCJKmainfont 控制，两者互不干扰，
%       因此英文显示为新罗马、中文保持规范字体。
% ============================================================

% —— 西文字体 ——
\setmainfont{Times New Roman}     % 英文正文衬线 = Times New Roman
\setsansfont{Arial}               % 西文无衬线（标题/图表标签可选）
\setmonofont{Courier New}         % 西文等宽（代码、算法）

% —— 辅助命令：在中文（宋体/黑体）环境中强制西文/数字用 Times New Roman ——
% 例：\tnr{2020110101} 用于封面"学号"等需要 TNR 阿拉伯数字处。
\newcommand{\tnr}[1]{{\fontspec{Times New Roman}#1}}

% —— 如需手动指定中文字体（例如 Linux 无系统字体时）——
% 取消下面注释并改为本机可用字体名，同时把 main.tex 的
% fontset=auto 改为 fontset=none：
% \setCJKmainfont{SimSun}          % 宋体（正文）
% \setCJKsansfont{SimHei}          % 黑体（无衬线/部分标题）
% \setCJKfamilyfont{zhkai}{KaiTi}  % 楷体
% \setCJKfamilyfont{zhfs}{FangSong} % 仿宋
